% !TeX spellcheck = pt_BR
\documentclass[12pt,portuguese]{article}

%----------------------------------------------------------------------
% PACOTES ESSENCIAIS E TRADUÇÃO
%----------------------------------------------------------------------
\usepackage[brazilian]{babel}     % Suporte à língua portuguesa
\usepackage[utf8]{inputenc}       % Codificação de entrada
\usepackage[T1]{fontenc}          % Codificação da fonte
\usepackage{lmodern}              % Fonte escalável

%----------------------------------------------------------------------
% PACOTES MATEMÁTICOS
%----------------------------------------------------------------------
\usepackage{amsmath, amsthm, amssymb, amstext, mathtools}
\usepackage{mathrsfs}             % Para fontes como \mathscr
\usepackage{dsfont}               % Para o indicador de conjunto \mathds{1}
\usepackage[all]{xy}              % Para diagramas comutativos
\usepackage{imakeidx}             % Criação de Lista de Símbolos / Índice

% Desigualdades sempre com símbolos inclinados
\let\leq\leqslant
\let\le\leqslant
\let\ge\geqslant
\let\geq\geqslant

%----------------------------------------------------------------------
% FORMATAÇÃO, LAYOUT E CORES
%----------------------------------------------------------------------
\usepackage{geometry}
\usepackage{varwidth}
\usepackage{graphicx}
\usepackage[pdftex]{xcolor}
\usepackage{titlesec}
\usepackage{appendix}
\usepackage{enumitem}
\usepackage{float}
\usepackage{booktabs}
\usepackage{array}
\usepackage{colortbl}
\usepackage{fontawesome5}
\usepackage{lucide-icons}
\usepackage{caption}
\usepackage{stackengine}
\captionsetup{width=0.8\textwidth}

% Definição de Cores
\definecolor{teal}{RGB}{15,118,110}
\definecolor{text}{RGB}{28,25,23}
\definecolor{muted}{RGB}{87,83,78}
\definecolor{card}{RGB}{255,255,255}
\definecolor{border}{RGB}{231,229,228}

% Definição de margens e espaçamento
\setlength{\topmargin}{-0.7in}
\setlength{\oddsidemargin}{0in}
\setlength{\textheight}{9.55in}
\setlength{\textwidth}{6.5in}
\setlength{\baselineskip}{12mm}

% Formatação dos títulos de seção
\titleformat{\section}{\Large\bfseries}{\thesection.}{0.5em}{}
\titleformat{\subsection}{\large\bfseries}{\thesubsection.}{0.5em}{}
\titleformat{\subsubsection}{\normalfont\itshape}{\thesubsubsection.}{0.6em}{}

%----------------------------------------------------------------------
% TIKZ E GRÁFICOS
%----------------------------------------------------------------------
\usepackage{tikz}
\usepackage{pgfplots}
\pgfplotsset{compat=1.17}
\usetikzlibrary{positioning,calc,fit,nfold,decorations.pathreplacing,calligraphy,tikzmark,patterns,arrows.meta,bending,hobby}

%----------------------------------------------------------------------
% REFERÊNCIAS, LINKS E METADADOS DO PDF
%----------------------------------------------------------------------
\usepackage[pdftex, colorlinks, linkcolor={red!50!black}, 
citecolor={red}, urlcolor={blue!80!black}]{hyperref}

\hypersetup{
pdftitle={Somas de Variáveis Aleatórias Contínuas e Mudança de Variáveis},
pdfauthor={L. Cioletti},
pdfsubject={Somas de variáveis aleatórias contínuas, fórmula da convolução, densidade da soma, mudança de variáveis, determinante jacobiano e densidade conjunta transformada},
pdfkeywords={somas de variáveis aleatórias contínuas, fórmula da convolução, densidade da soma, mudança de variáveis, determinante jacobiano, densidade conjunta, MAT-UnB}
}

% --- PACOTE E COMANDO PARA A LICENÇA CREATIVE COMMONS ---
\usepackage[type={CC},modifier={by-nc-sa},version={4.0}]{doclicense}
\newcommand\blfootnote[1]{%
\begingroup
\renewcommand\thefootnote{}\footnote{#1}%
\addtocounter{footnote}{-1}%
\endgroup
}

%----------------------------------------------------------------------
% ALIASCNT E AMBIENTES DE TEOREMA (Cleveref)
%----------------------------------------------------------------------
\usepackage{aliascnt}
\theoremstyle{definition}
\newtheorem{theorem}{Teorema}

\newaliascnt{definition}{theorem}
\newtheorem{definition}[definition]{Definição}
\aliascntresetthe{definition}

\newaliascnt{example}{theorem}
\newtheorem{example}[example]{Exemplo}
\aliascntresetthe{example}

\newaliascnt{proposition}{theorem}
\newtheorem{proposition}[proposition]{Proposição}
\aliascntresetthe{proposition}

\newaliascnt{corollary}{theorem}
\newtheorem{corollary}[corollary]{Corolário}
\aliascntresetthe{corollary}

\newaliascnt{lemma}{theorem}
\newtheorem{lemma}[lemma]{Lema}
\aliascntresetthe{lemma}

\newaliascnt{remark}{theorem}
\newtheorem{remark}[remark]{Observação}
\aliascntresetthe{remark}

\newtheoremstyle{plainexercise}
  {3pt}{3pt}{\normalfont}{}{\bfseries}{.}{.5em}
  {\thmnumber{#2}\thmnote{ #3}}
\theoremstyle{plainexercise}
\newtheorem{exercise}{Exercício}

% Versão marcada dos exercícios: a estrela fica pendurada na margem esquerda,
% como nas notas anteriores.
\newtheoremstyle{starredexercise}
  {3pt}{3pt}{\normalfont}{}{\bfseries}{.}{.5em}
  {\llap{\textasteriskcentered\,}\thmnumber{#2}\thmnote{ #3}}
\theoremstyle{starredexercise}
\newtheorem{starredexercise}[exercise]{Exercício}
\theoremstyle{definition}

% Ambientes não numerados usados no texto
\newtheorem*{property}{Propriedade}
\newtheorem*{solution}{Solução}

% Se um enumerate for o primeiro conteúdo de um teorema/exercício (cabeçalho
% ainda pendente, \@inlabel true), o \leavevmode dispara o \everypar do amsthm
% que coloca o cabeçalho na própria linha, antes da lista iniciar.
\makeatletter
\setlist[enumerate]{before={\if@inlabel\leavevmode\fi}}
\makeatother

% Nomes em português para os ambientes
\AtBeginDocument{\renewcommand{\proofname}{\normalfont\textbf{Prova}}}
\renewcommand{\qedsymbol}{$\blacksquare$}

% Cleveref (deve vir depois de hyperref e amsthm)
\usepackage[brazilian,nameinlink]{cleveref}

\crefname{theorem}{teorema}{teoremas}
\Crefname{theorem}{Teorema}{Teoremas}
\crefname{definition}{definição}{definições}
\Crefname{definition}{Definição}{Definições}
\crefname{lemma}{lema}{lemas}
\Crefname{lemma}{Lema}{Lemas}
\crefname{proposition}{proposição}{proposições}
\Crefname{proposition}{Proposição}{Proposições}
\crefname{corollary}{corolário}{corolários}
\Crefname{corollary}{Corolário}{Corolários}
\crefname{example}{exemplo}{exemplos}
\Crefname{example}{Exemplo}{Exemplos}
\crefname{exercise}{exercício}{exercícios}
\Crefname{exercise}{Exercício}{Exercícios}
\crefname{remark}{observação}{observações}
\Crefname{remark}{Observação}{Observações}
\crefname{section}{seção}{seções}
\Crefname{section}{Seção}{Seções}

%----------------------------------------------------------------------
% MACROS ADICIONAIS
%----------------------------------------------------------------------
\let\oldemptyset\emptyset
\let\emptyset\varnothing
\graphicspath{{./}{../arquivos-aux/}}

% Seno na convenção em português
\renewcommand{\sin}{\operatorname{sen}}

%----------------------------------------------------------------------
% DOCUMENTO PRINCIPAL
%----------------------------------------------------------------------
\begin{document}
% --- CABEÇALHO COM O LOGO DA UnB ---
\vspace*{-0.8cm}
\begin{figure}[H]
\includegraphics{logo-mat-unb.pdf}
\end{figure}


% --- TÍTULO, AUTOR E DATA ---
\vspace*{0.3cm}
\begin{center}
{\LARGE\bf Somas de Variáveis Aleatórias Contínuas \\[0.2cm] e Mudança de Variáveis}
\vspace{1cm}

{\large \bf L. Cioletti}
\vspace{0.3cm}

{\footnotesize Outubro de 2026}
\end{center}

% --- RESUMO (ABSTRACT) ---
\begin{abstract}
Nestas notas, damos continuidade ao estudo dos vetores aleatórios contínuos.
Começamos determinando a densidade da soma $Z = X + Y$ de duas variáveis
aleatórias conjuntamente contínuas e, no caso de independência, obtemos a
fórmula da convolução. Como aplicação, mostramos que a soma de duas variáveis
independentes com distribuição gama de mesma taxa tem distribuição gama com os
parâmetros de forma somados, o que recupera, por outro caminho, a relação entre
as funções beta e gama. Em seguida, estudamos a fórmula do jacobiano para a
distribuição conjunta de um par $(U, V) = (u(X, Y), v(X, Y))$, que permite
transferir a densidade conjunta de $(X, Y)$ para o novo par por meio do
determinante jacobiano da transformação inversa. Ilustramos o método com uma
transformação que produz uma variável uniforme independente de uma variável
com distribuição gama. A exposição segue de perto as Seções 6.4 e 6.5 do
Capítulo 6 de \cite{GrimmettWelsh2014}.
\end{abstract}

% --- LICENÇA NO RODAPÉ ---
\blfootnote{%
\doclicenseImage[imagewidth=2.6em]\ 
\begingroup
\hypersetup{urlcolor=darkgray}
\href{\doclicenseURL}{%
\doclicenseName - Attribution-NonCommercial-ShareAlike 4.0 International%
}%
\endgroup
}

%----------------------------------------------------------------------
% CONTEÚDO
%----------------------------------------------------------------------


\section{Somas de Variáveis Aleatórias Contínuas}\label{sec:somas-continuas}
Nas notas de aula anteriores sobre vetores aleatórios contínuos, aprendemos como
descrever a distribuição conjunta de um par $(X, Y)$ absolutamente contínuo 
por meio de uma função densidade conjunta $f_{X,Y}$ e vimos como recuperar as distribuições de $X$ e
de $Y$ integrando a densidade conjunta em relação à outra coordenada. Também
caracterizamos a independência pela fatoração da densidade conjunta, quase em
todo ponto, como o produto das densidades marginais. Nestas notas, utilizamos
essas ferramentas para responder a uma pergunta natural: dada uma função das
coordenadas, como a soma $Z = X + Y$, qual é a sua distribuição? Mais adiante,
trataremos do caso geral em que o par de interesse é obtido de $(X, Y)$ por uma
transformação $(U, V) = (u(X, Y), v(X, Y))$, e para isso desenvolveremos a
fórmula do jacobiano.

A soma de variáveis aleatórias aparece em muitas aplicações e merece um estudo próprio. 
Para determinar sua densidade, começaremos calculando a função de distribuição da soma. 
Uma mudança de variáveis permitirá escrever essa função como uma integral cuja 
derivada fornece a densidade procurada. 
Posteriormente, usaremos mudanças de variáveis para determinar a densidade 
conjunta de pares obtidos por transformações que podem depender de formas mais gerais do
vetor aleatório $(X,Y)$. 
 

\subsection{A Densidade da Soma}\label{subsec:densidade-soma}
Como a densidade de uma variável aleatória é, quando existe, a derivada de sua
função de distribuição, calculamos primeiro $\mathbb{P}(Z \leqslant z)$ e em
seguida derivamos em relação a $z$. Logo se $Z\equiv X+Y$, então o evento $\{Z \leqslant z\}$ é o mesmo que
$\{X + Y \leqslant z\}$, ou seja, é o evento correspondente ao vetor aleatório $(X, Y)$
pertencer à região
\[
  A = \{(x, y) \in \mathbb{R}^2 : x + y \leqslant z\}.
\]

Supondo que o vetor aleatório $(X,Y)$ possua uma função de densidade conjunta
segue do Teorema da Mudança de Variáveis que a probabilidade de $(X, Y)$ pertencer a $A$ é a
integral da densidade conjunta sobre essa região. Logo
\begin{align*}
  \mathbb{P}(Z \leqslant z)
  &= \mathbb{P}(X + Y \leqslant z)
  \\[0.3cm]
  &= \iint_A f_{X,Y}(x, y)\,dy\,dx.
\end{align*}

Para calcular a integral, descrevemos $A$ por meio das desigualdades
$-\infty < x < \infty$ e $-\infty < y \leqslant z - x$. 
\begin{figure}[H]
\centering
\begin{tikzpicture}[scale=0.9, >=Stealth]
  % Janela de visualização do semiplano: a região não é limitada.
  \fill[teal!15]
    (-2.6,-2.5) -- (4.9,-2.5) -- (-2.3,4.7) -- (-2.6,4.7) -- cycle;

  \draw[->, thick] (-2.9,0) -- (5.3,0) node[below] {$x$};
  \draw[->, thick] (0,-2.7) -- (0,5.0) node[left] {$y$};
  % A fronteira pertence a A, por isso a reta é contínua.
  \draw[thick, muted] (-2.3,4.7) -- (4.9,-2.5);
  \node[rotate=-45, text=text, above=3pt] at (3.90,-1.50)
    {$x+y=z$};
  \node[teal!80!black] at (-2,-2) {$A$};
  \node[teal!80!black] at (2.95,-1.95) {$y\leqslant z-x$};

  % Interceptos, desenhados para z>0 apenas para facilitar a leitura.
  \fill[muted] (2.4,0) circle (1.5pt)
    node[above=7pt, right] {$z$};
  \fill[muted] (0,2.4) circle (1.5pt) node[above=5pt, right] {$z$};
  \node[below left=2pt] at (0,0) {$0$};

  % Para x=x_0, a seção vertical de A é uma semirreta fechada.
  \draw[red!75, dashed] (0,1.05) -- (1.35,1.05);
  \node[left=3pt] at (0,1.05) {$z-x_0$};
  \draw[red!75, thick] (1.35,1.05) -- (1.35,-2.4);
  \fill[red!75] (1.35,1.05) circle (1.8pt);
  \fill[red!75] (1.35,0) circle (1.8pt);
  \node[below left] at (1.35,0) {$x_0$};
  \node[red!75, below=3pt] at (1.35,-2.4) {$-\infty$};
\end{tikzpicture}
\caption{O semiplano $A=\{(x,y):x+y\leqslant z\}$, sombreado,
inclui a reta $x+y=z$ e se estende além da janela desenhada. Para cada
$x=x_0$, a seção vertical indicada em vermelho é dada por
$-\infty<y\leqslant z-x_0$.}
\label{fig:semiplano-densidade-soma}
\end{figure}
Escrevendo os limites
de integração de forma explícita, obtemos
\[
  \mathbb{P}(Z \leqslant z)
  = \int_{-\infty}^{\infty} \int_{-\infty}^{z-x} f_{X,Y}(x, y)\,dy\,dx.
\]

Para cada $x$ fixado, fazemos na integral interna a substituição $v=x+y$.
Como $y=v-x$ e $dy=dv$, obtemos
\[
  \mathbb{P}(Z \leqslant z)
  = \int_{-\infty}^{\infty} \left[ \int_{-\infty}^{z} f_{X,Y}(x, v-x)\,dv \right] dx.
\]
Agora o limite superior da integral interna não depende de $x$. Como o
integrando é não negativo e mensurável, o Teorema de Tonelli permite trocar a
ordem de integração:
\[
  \mathbb{P}(Z \leqslant z)
  = \int_{-\infty}^{z} \left[ \int_{-\infty}^{\infty} f_{X,Y}(x, v-x)\,dx \right] dv.
\]

Próximo passo e calcular a derivada da expressão acima, em relação a $z$. 
Para deixar mais claro o argumento, denotamos a integrando mais interno da integral
dupla acima por
\[
  H(v) = \int_{-\infty}^{\infty} f_{X,Y}(x, v-x)\,dx.
\]
Logo, temos da igualdade anterior que
\begin{align}\label{eq-aux1-dens-soma}
  \mathbb{P}(Z \leqslant z) = \int_{-\infty}^{z} H(v)\,dv.
\end{align}
A função $H$ é Borel-mensurável e não negativa, por ser a integral de uma
função com essas propriedades. Além disso, quando $z$ tende a $+\infty$, o lado
esquerdo de \eqref{eq-aux1-dens-soma} tende a $1$, e portanto $\int_{-\infty}^{\infty} H(v)\,dv = 1 <
\infty$. Assim, $H$ é integrável. Pelo teorema fundamental do cálculo para a
integral de Lebesgue, a derivada da função $z \mapsto \int_{-\infty}^{z}
H(v)\,dv$ é igual a $H(z)$ em quase todo ponto. Assim concluímos que uma densidade de
$Z$ é dada por
\begin{equation}\label{eq:densidade-soma}
  f_Z(z) = \int_{-\infty}^{\infty} f_{X,Y}(x, z-x)\,dx
  \quad \text{para quase todo ponto } z \in \mathbb{R}.
\end{equation}
Observamos que sem hipóteses adicionais
sobre a regularidade de $f_{X,Y}$ igualdade obtida acima é válida apenas 
para quase todo ponto. 
A razão é que só é possível garantir a existência da derivada do lado direito da
expressão \eqref{eq-aux1-dens-soma}, com relação a $z$, em quase todo ponto. 
Entretanto, se $f_{X,Y}$ for uma função contínua então podemos afirmar que a derivada
existe em todo ponto. 
De qualquer maneira, como os valores de uma densidade em um conjunto 
de medida nula não alteram as probabilidades, não é necessário
assumir que $f_{X,Y}$ em todo ponto.

\subsection{O Caso Independente e a Fórmula da Convolução}\label{subsec:formula-convolucao}
Um caso particular importante é quando $X$ e $Y$ são variáveis aleatórias independentes. 
Neste caso, a densidade conjunta se fatora, em quase todo ponto, como $f_{X,Y}(x, y) =
f_X(x)\,f_Y(y)$, e a substituição em \eqref{eq:densidade-soma} fornece a
fórmula da convolução.

\begin{theorem}[Fórmula da Convolução]\label{thm:formula-convolucao}
  Sejam $X$ e $Y$ variáveis aleatórias independentes e absolutamente contínuas,
  com funções densidade $f_X$ e $f_Y$. Então $X$ e $Y$ são conjuntamente
  absolutamente contínuas e a variável aleatória $Z = X + Y$ tem densidade
  dada por
  \begin{equation}\label{eq:convolucao}
    f_Z(z) = \int_{-\infty}^{\infty} f_X(x)\,f_Y(z - x)\,dx
    \quad \text{para quase todo ponto } z \in \mathbb{R}.
  \end{equation}
\end{theorem}

\begin{proof}
  Como $X$ e $Y$ são independentes e têm densidades $f_X$ e $f_Y$, o
  produto $f_X(x)\,f_Y(y)$ é uma densidade conjunta do par $(X, Y)$, conforme
  mostramos nas notas sobre vetores contínuos. Em particular, a densidade
  conjunta satisfaz $f_{X,Y}(x, y) = f_X(x)\,f_Y(y)$ quase em todo ponto.
  Substituindo essa fatoração em \eqref{eq:densidade-soma}, obtemos
  \eqref{eq:convolucao}. A
  alteração da densidade conjunta em um conjunto de área nula não modifica o
  valor da integral, o que justifica a substituição.
\end{proof}

Na linguagem da análise, a equação \eqref{eq:convolucao} diz que $f_Z$ é a
\emph{convolução} das funções $f_X$ e $f_Y$, e escrevemos $f_Z = f_X * f_Y$.
De modo geral, a convolução de duas funções integráveis $f$ e $g$ é definida
por
\[
  (f * g)(z) = \int_{-\infty}^{\infty} f(x)\,g(z - x)\,dx,
\]
sempre que a integral fizer sentido. Dizer que a densidade da soma de variáveis
aleatórias independentes é a convolução de suas densidades é, portanto, o mesmo
que dizer que a distribuição de $X + Y$ é determinada pela operação acima.
Voltaremos a esse ponto nos exemplos a seguir, nos quais aplicamos a
\Cref{thm:formula-convolucao} a pares de variáveis aleatórias concretos.

\subsection{Somas de Variáveis Gama}\label{subsec:somas-gama}
\begin{example}[Soma de Variáveis Gama Independentes]\label{ex:soma-gamas}
  Sejam $X$ e $Y$ variáveis aleatórias independentes, com
  $X\sim\mathrm{Gama}(s,\lambda)$ e $Y\sim\mathrm{Gama}(t,\lambda)$, onde
  $s,t,\lambda>0$. Recordamos, das notas sobre funções de distribuição, que a
  densidade gama de parâmetros $w>0$ e $\lambda>0$ é dada por
  \[
    f(x) =
    \begin{cases}
      \dfrac{1}{\Gamma(w)}\,\lambda^w x^{w-1} e^{-\lambda x}, & \text{se } x > 0, \\[0.3cm]
      0, & \text{se } x \leqslant 0,
    \end{cases}
  \]
  onde $w$ é o parâmetro de forma e $\lambda$ é a taxa. Escrevemos $Z=X+Y$ e
  vamos mostrar que $Z$ também tem distribuição gama, com os parâmetros de
  forma combinados de modo natural.

  Como $X$ e $Y$ são independentes e cada uma delas tem densidade, segue da
  fórmula da convolução, \eqref{eq:convolucao}, que a densidade de $Z$ é
  \begin{equation}
    f_Z(z) = \int_{-\infty}^{\infty} f_X(x) f_Y(z-x)\,dx.
    \label{eq:t2-convolucao}
  \end{equation}
  O integrando $f_X(x)f_Y(z-x)$ é não negativo, e cada fator tem suporte
  restrito: $f_X(x)$ se anula para $x\leqslant 0$, enquanto $f_Y(z-x)$ se
  anula quando $z-x\leqslant 0$, isto é, quando $x\geqslant z$. Assim, para
  $z$ fixado, o integrando é nulo fora do intervalo $0<x<z$. Quando
  $z>0$ esse intervalo é não vazio e a integral se reduz a ele; quando
  $z\leqslant 0$ o intervalo é vazio e o integrando é nulo em toda a reta.
  Portanto
  \[
    f_Z(z) =
    \begin{cases}
      \displaystyle\int_0^z f_X(x) f_Y(z-x)\,dx, & \text{se } z > 0, \\[0.3cm]
      0, & \text{se } z \leqslant 0.
    \end{cases}
  \]

  Resta calcular a integral no caso $z>0$. Substituindo as duas densidades,
  \begin{align*}
    f_Z(z)
    &= \int_0^z \frac{1}{\Gamma(s)}\lambda^s x^{s-1} e^{-\lambda x}
       \cdot \frac{1}{\Gamma(t)}\lambda^t (z-x)^{t-1} e^{-\lambda(z-x)}\,dx
    \\[0.3cm]
    &= \frac{\lambda^{s+t}}{\Gamma(s)\Gamma(t)}
       \int_0^z x^{s-1}(z-x)^{t-1} e^{-\lambda x - \lambda(z-x)}\,dx
    \\[0.3cm]
    &= \frac{\lambda^{s+t}}{\Gamma(s)\Gamma(t)}\, e^{-\lambda z}
       \int_0^z x^{s-1}(z-x)^{t-1}\,dx,
  \end{align*}
  pois $e^{-\lambda x}e^{-\lambda(z-x)}=e^{-\lambda z}$ e o fator
  $e^{-\lambda z}$ não depende de $x$, podendo ser retirado da integral.

  Para calcular a integral que permanece, fazemos a substituição $y=x/z$, ou
  seja, $x=zy$, com $dx=z\,dy$. Enquanto $x$ percorre o intervalo $(0,z)$, a
  nova variável $y$ percorre $(0,1)$, e
  \[
    x^{s-1}(z-x)^{t-1}
    = (zy)^{s-1}\bigl(z(1-y)\bigr)^{t-1}
    = z^{s+t-2} y^{s-1}(1-y)^{t-1}.
  \]
  Multiplicando ainda pelo fator $z$ proveniente do diferencial, obtemos
  \begin{equation}
    \int_0^z x^{s-1}(z-x)^{t-1}\,dx
    = z^{s+t-1}\int_0^1 y^{s-1}(1-y)^{t-1}\,dy
    = z^{s+t-1} B(s,t),
    \label{eq:t2-integral-beta}
  \end{equation}
  onde reconhecemos na última integral a função beta $B(s,t)$, definida nas
  notas sobre funções de distribuição.

  Reunindo os dois últimos passos, concluímos que, para $z>0$,
  \begin{equation}
    f_Z(z) = C\, z^{s+t-1} e^{-\lambda z},
    \label{eq:densidade-soma-gamas}
  \end{equation}
  em que a constante $C$ é dada por
  \[
    C = \frac{\lambda^{s+t}}{\Gamma(s)\Gamma(t)}\,B(s,t).
  \]

  Para identificar $C$, usamos a relação beta--gama demonstrada nas notas
  sobre funções de distribuição, segundo a qual
  \[
    B(s,t) = \frac{\Gamma(s)\Gamma(t)}{\Gamma(s+t)}
  \]
  para quaisquer $s,t>0$. Substituindo essa expressão em $C$, os fatores
  $\Gamma(s)$ e $\Gamma(t)$ se cancelam e obtemos
  \[
    C = \frac{\lambda^{s+t}}{\Gamma(s+t)}.
  \]

  Portanto
  \[
    f_Z(z) =
    \begin{cases}
      \dfrac{1}{\Gamma(s+t)}\,\lambda^{s+t} z^{s+t-1} e^{-\lambda z}, & \text{se } z > 0, \\[0.3cm]
      0, & \text{se } z \leqslant 0.
    \end{cases}
  \]
  Comparando com a densidade gama, reconhecemos em $f_Z$ a densidade de
  parâmetros $s+t$ e $\lambda$. Segue que
  \[
    X+Y \sim \mathrm{Gama}(s+t,\lambda).
  \]
  Em palavras: a soma de duas variáveis independentes com distribuição gama de
  mesma taxa tem distribuição gama cujos parâmetros de forma se somam, com a
  taxa preservada.

  O cálculo acima também fornece uma maneira indireta de calcular $B(s,t)$. 
  Vamos imaginar que não conhecemos a
  relação beta--gama. Poderíamos, neste caso, determinar $C$ exigindo que a densidade em
  \eqref{eq:densidade-soma-gamas} tenha integral $1$. Isto força a seguinte igualdade
  $C=\lambda^{s+t}/\Gamma(s+t)$. Comparando esse valor com a expressão
  $C=\lambda^{s+t}B(s,t)/(\Gamma(s)\Gamma(t))$ obtida acima, poderíamos concluir que
  \[
    B(s,t) = \frac{\Gamma(s)\Gamma(t)}{\Gamma(s+t)},
  \]
  isto é, o argumento de convolução reobtém a relação beta--gama.
\end{example}

\begin{exercise}\label{exercise:densidade-soma-exercicio}
  As variáveis aleatórias $X$ e $Y$ têm função densidade conjunta dada por
  \[
    f_{X,Y}(x, y) =
    \begin{cases}
      \frac{1}{2}(x + y)e^{-x-y}, & \text{se } x > 0 \text{ e } y > 0, \\[0.3cm]
      0, & \text{caso contrário}.
    \end{cases}
  \]
  Determine a função densidade de $X + Y$.
\end{exercise}

\begin{exercise}\label{exercise:soma-qui-quadrado}
  Sejam $X$ e $Y$ variáveis aleatórias independentes com distribuição
  qui-quadrado com $m$ e $n$ graus de liberdade, respectivamente.
  Mostre que $X + Y$ tem distribuição qui-quadrado com $m + n$ graus de
  liberdade.

  Recorde, das notas sobre funções de distribuição, que a distribuição
  qui-quadrado com $n$ graus de liberdade é a distribuição gama com
  parâmetros de forma $n/2$ e de taxa $1/2$.
\end{exercise}

\begin{exercise}\label{exercise:soma-normais}
  Sejam $X$ e $Y$ variáveis aleatórias independentes, cada uma com
  distribuição normal padrão, isto é, com média $0$ e variância $1$.
  Determine a distribuição de $X + Y$.
\end{exercise}

\section{Mudança de Variáveis}\label{sec:mudanca-variaveis}
A seguinte situação ocorre com frequência na Teoria de Probabilidade. Sejam $X$ e $Y$ variáveis
aleatórias e sejam $u, v : \mathbb{R}^2 \to \mathbb{R}$ funções. O que podemos
dizer sobre a distribuição conjunta do vetor $(U, V)$ de variáveis aleatórias
definido por
\[
  U = u(X, Y), \qquad V = v(X, Y)?
\]
Esta é a pergunta que orienta a presente seção. Vamos apresentar uma resposta
para o caso em que $X$ e $Y$ são conjuntamente absolutamente contínuas e as funções $u$ e $v$
satisfazem certas condições que permitem recorrer à teoria usual de mudança de
variáveis dentro de uma integral.

\subsection{O Jacobiano}\label{subsec:jacobiano}
Seja $T$ a aplicação de $\mathbb{R}^2$ em $\mathbb{R}^2$ dada por
$T(x, y) = (u, v)$, onde $u = u(x, y)$ e $v = v(x, y)$, e suponhamos que $T$
seja uma bijeção entre um domínio $D \subseteq \mathbb{R}^2$ e um contradomínio
$S \subseteq \mathbb{R}^2$. 
\begin{figure}[H]
\centering
\begin{tikzpicture}[scale=1.05, >=Stealth]
  % Plano de partida e representação abstrata do domínio.
  \fill[teal!15] (0,0) circle (1.4);
  \draw[teal, line width=0.9pt] (0,0) circle (1.4);
  \draw[->, thick] (-2.22,0) -- (2.34,0) node[below] {$x$};
  \draw[->, thick] (0,-2.22) -- (0,2.34) node[left] {$y$};
  \node[teal!80!black] at (-0.65,0.65) {$D$};

  % Plano de chegada e representação abstrata do contradomínio.
  % Os eixos horizontais dos dois planos ficam na mesma altura.
  \begin{scope}[xshift=5.2cm]
    \fill[teal!15] (0,0) rectangle (2.35,1.8);
    \draw[teal, line width=0.9pt] (0,0) rectangle (2.35,1.8);
    \draw[->, thick] (-0.42,0) -- (3.54,0) node[below] {$u$};
    \draw[->, thick] (0,-0.42) -- (0,2.82) node[left] {$v$};
    \node[teal!80!black] at (1.175,0.9) {$S$};
  \end{scope}

  % A seta indica a aplicação entre as regiões, sem fixar uma fórmula.
  \draw[->, teal!80!black, line width=0.9pt]
    (1.3,1.2) to[out=30, in=150, looseness=1.2]  node[pos=0.5, above=4pt] {$T$} (4.6,1.0);
  % Reflexão da curva no eixo horizontal, com o sentido invertido.
  \draw[->, teal!80!black, line width=0.9pt]
    (4.6,-1.0) to[out=-150, in=-30, looseness=1.2]
    node[pos=0.5, below=4pt] {$T^{-1}$} (1.3,-1.2);
\end{tikzpicture}
\caption{Representação esquemática de uma bijeção $T:D\to S$
entre duas regiões do plano e de sua inversa $T^{-1}:S\to D$.}
\label{fig:bijecao-mudanca-variaveis}
\end{figure}
Nesse caso, podemos inverter $T$ e obter uma
bijeção $T^{-1} : S \to D$. Isto é, para cada $(u, v) \in S$ existe um único
ponto $(x, y) = T^{-1}(u, v)$ em $D$, e escrevemos
\[
  x = x(u, v), \qquad y = y(u, v).
\]

O jacobiano de $T^{-1}$ é definido como o determinante
\begin{equation}
\label{eq:jacobiano}
  J(u, v)
  = \left|
    \begin{array}{>{\displaystyle}c@{\hspace{0.8cm}}>{\displaystyle}c}
      \frac{\partial x}{\partial u} & \frac{\partial x}{\partial v} 
      \\[0.6cm]
      \frac{\partial y}{\partial u} & \frac{\partial y}{\partial v}
    \end{array}
    \right|
  = \frac{\partial x}{\partial u} \frac{\partial y}{\partial v}
  - \frac{\partial y}{\partial u} \frac{\partial x}{\partial v},
\end{equation}
e supomos que essas derivadas parciais existem e são contínuas em todos os
pontos de $S$.

A teoria padrão das integrais múltiplas nos diz como efetuar uma mudança de
variáveis dentro da integral. 
Se $g : D \to \mathbb{R}$ é uma função integrável, então
\begin{equation}
\label{eq:mudanca-variaveis-integral}
  \iint_D g(x, y)\,dx\,dy
  = \iint_{T(D)} g\bigl(x(u, v), y(u, v)\bigr)\,|J(u, v)|\, dudv
\end{equation}
ou equivalentemente,
\begin{equation}
\label{eq:mudanca-variaveis-integral-v2}
  \iint_{T^{-1}(S)} g(x, y)\,dx\,dy
  = \iint_{S} g\bigl(x(u, v), y(u, v)\bigr)\,|J(u, v)|\, dudv
\end{equation}


A presença do fator $|J(u, v)|$ nessa fórmula se explica da seguinte maneira:
o valor absoluto do determinante jacobiano mede a dilatação local de área
produzida pela aplicação inversa $T^{-1}$. Em outras palavras, um retângulo
infinitesimal de área $du\,dv$ em torno do ponto $(u, v) \in S$ é levado por
$T^{-1}$ a uma região infinitesimal em torno de $(x, y) = T^{-1}(u, v)$ cuja
área é $|J(u, v)|\,du\,dv$. Esse é o motivo pelo qual a integral sobre $A$ se
transforma na integral sobre $T(A)$ mediante a multiplicação do integrando pelo
fator $|J(u, v)|$.
\begin{figure}[H]
\centering
\begin{tikzpicture}[scale=1.05, >=Stealth]
  % Plano x-y: imagem do pequeno retângulo pela aplicação inversa.
  \draw[->, thick] (-2.22,0) -- (2.34,0) node[below] {$x$};
  \draw[->, thick] (0,-2.22) -- (0,2.34) node[left] {$y$};
  % Curvaturas suaves sugerem que T^{-1} pode ser não linear.
  \filldraw[fill=teal!15, draw=teal!80!black, line width=0.9pt]
    (0.524,0.516)
    to[out=2, in=202] (1.42,0.708)
    to[out=56, in=256] (1.676,1.284)
    to[out=182, in=22] (0.78,1.092)
    to[out=236, in=76] (0.524,0.516) -- cycle;
  \fill[teal!80!black] (1.1,0.9) circle (1.8pt);
  \node[below=4pt] at (1.4,0.74) {$(x,y)$};
  \node at (0,-2.65) {$\text{Área}\approx |J(u,v)|\,du\,dv$};

  % Plano u-v: o retângulo infinitesimal e seus dois comprimentos.
  \begin{scope}[xshift=5.2cm]
    \draw[->, thick] (-0.42,0) -- (3.54,0) node[below] {$u$};
    \draw[->, thick] (0,-0.42) -- (0,2.82) node[left] {$v$};
    \filldraw[fill=teal!15, draw=teal!80!black, line width=0.9pt]
      (0.65,0.45) rectangle (1.7,1.25);
    \fill[teal!80!black] (1.175,0.85) circle (1.8pt);
    \node[above=4pt] at (1.175,1.3) {$(u,v)$};
    \draw[<->, teal!80!black] (0.65,0.25) -- (1.7,0.25)
      node[midway, below=2pt] {$du$};
    \draw[<->, teal!80!black] (1.9,0.45) -- (1.9,1.25)
      node[midway, right=2pt] {$dv$};
    \node at (1.175,-2.65) {$\text{Área}=du\,dv$};
  \end{scope}

  \draw[->, teal!80!black, line width=0.9pt]
    (5.5,1.0) to[out=165, in=15, looseness=1.2]
    node[pos=0.5, above=4pt] {$T^{-1}$} (1.9,1.0);
\end{tikzpicture}
\caption{Dilatação local de área pela aplicação inversa $T^{-1}$.
O quadrilátero de lados curvos representa a imagem do pequeno
retângulo; a razão entre as áreas tende a $|J(u,v)|$ quando o retângulo
encolhe em torno de $(u,v)$.}
\label{fig:jacobiano-area-local}
\end{figure}


No caso unidimensional, onde $T:[c,d]\to [a,b]$ é uma função
estritamente crescente e diferenciável, 
a fórmula \eqref{eq:mudanca-variaveis-integral-v2} se
reduz à conhecida fórmula de substituição simples
\[
  \int_{T^{-1}([a,b])} g(x)\, dx
  =
  \int_{T^{-1}(a)}^{T^{-1}(b)} g(t)\,dt 
  = 
  \int_{a}^{b} g(T^{-1}(u))\, (T^{-1})'(u)\,du,
\]
o que ajuda a fixar a interpretação do jacobiano como fator de escala local. A
seguir, vamos combinar \eqref{eq:mudanca-variaveis-integral} com a
probabilidade de regiões borelianas, estudada nas notas sobre vetores
contínuos, para transferir a densidade conjunta de $(X, Y)$ para a densidade
conjunta de $(U, V)$.

\subsection{A Fórmula do Jacobiano}\label{subsec:formula-jacobiano}
\begin{theorem}[Fórmula do Jacobiano]\label{thm:formula-jacobiano}
Sejam $X$ e $Y$ variáveis aleatórias conjuntamente contínuas, com função
densidade conjunta $f_{X,Y}$, e seja
$D = \{(x, y) : f_{X,Y}(x, y) > 0\}$.
Suponha que a aplicação $T:D\to S$
seja uma bijeção de $D$ sobre um conjunto $S \subseteq \mathbb{R}^2$,
dada por $T(x, y) = (u(x, y), v(x, y))$ e que
$T^{-1}$ tenha derivadas parciais contínuas em $S$, satisfeitas as
condições de regularidade que garantem a validade da fórmula de mudança de
variáveis \eqref{eq:mudanca-variaveis-integral} (por exemplo, quando o
jacobiano $J$ definido em \eqref{eq:jacobiano} não se anula em $S$). Então
o par $(U, V) = (u(X, Y), v(X, Y))$ é conjuntamente contínuo e sua função
densidade conjunta é dada por
\begin{equation}\label{eq:densidade-transformada}
  f_{U,V}(u, v) =
  \begin{cases}
    f_{X,Y}\bigl(x(u, v), y(u, v)\bigr)\,|J(u, v)|
    , & \text{se } (u, v) \in S,
    \\[0.2cm]
    0, & \text{caso contrário}.
  \end{cases}
\end{equation}
\end{theorem}

\begin{proof}
Vamos mostrar que a função $h$ definida por
\[
  h(u, v) =
  \begin{cases}
    f_{X,Y}\bigl(x(u, v), y(u, v)\bigr)\,|J(u, v)|
    , & \text{se } (u, v) \in S,
    \\[0.2cm]
    0, & \text{caso contrário},
  \end{cases}
\]
é uma função densidade conjunta para $(U, V)$. Primeiro observamos que
$(X, Y)$ assume valores em $D$ com probabilidade um, pois fora de $D$ a
densidade $f_{X,Y}$ se anula. Como $T$ leva $D$ sobre $S$, segue que
$(U, V) = T(X, Y)$ assume valores em $S$ com probabilidade um, de modo que
não há massa de probabilidade de $(U, V)$ fora de $S$. Além disso, $h$ é
não negativa, pois $f_{X,Y} \geqslant 0$ e $|J| \geqslant 0$, e é
Borel-mensurável, pois $f_{X,Y}$ é Borel-mensurável e $x$, $y$ e $J$ são
contínuas em $S$. Resta verificar que $h$ reproduz, por integração, as
probabilidades dos eventos que envolvem $(U, V)$.

Seja $B \subseteq S$ um boreliano e defina $A = T^{-1}(B)$, de modo que
$A \subseteq D$ e $T(A) = B$. Como $T^{-1}$ é contínua segue que $A$ é um boreliano
de $\mathbb{R}^2$. Vamos relacionar as probabilidades dos dois eventos por
meio da bijeção. Para todo $\omega$ no espaço amostral temos 
\[
  (X(\omega), Y(\omega)) \in A
  \quad \Longleftrightarrow \quad
  T\bigl(X(\omega), Y(\omega)\bigr) \in T(A) = B,
\]
isto é, $(X, Y) \in A$ se, e somente se, $(U, V) \in B$. Logo,
\begin{equation}\label{eq:probabilidade-bijecao}
  \mathbb{P}\bigl((U, V) \in B\bigr)
  = \mathbb{P}\bigl((X, Y) \in A\bigr).
\end{equation}
Já que o vetor $(X,Y)$ possui densidade conjunta podemos afirmar que o lado direito da igualdade acima é 
dado por
\[
  \mathbb{P}\bigl((X, Y) \in A\bigr)
  = \iint_A f_{X,Y}(x, y)\,dx\,dy.
\]
Agora aplicamos a fórmula de mudança de variáveis
\eqref{eq:mudanca-variaveis-integral} à integral que aparece no lado direito da igualdade acima, 
com a função integrável sendo $f_{X,Y}$ e $T(A) = B$. Como $J$ é o jacobiano da
inversa $T^{-1}$, temos
\[
  \iint_A f_{X,Y}(x, y)\,dx\,dy
  = \iint_B f_{X,Y}\bigl(x(u, v), y(u, v)\bigr)\,|J(u, v)|\,du\,dv
  = \iint_B h(u, v)\,du\,dv.
\]
Reunindo as duas igualdades, obtemos
\[
  \mathbb{P}\bigl((U, V) \in B\bigr) = \iint_B h(u, v)\,du\,dv
\]
para todo boreliano $B \subseteq S$. Essa é exatamente a propriedade que
caracteriza $h$ como densidade conjunta de $(U, V)$, e isso conclui a
prova de que o vetor aleatório $(U, V)$ possui densidade conjunta.
\end{proof}

A identificação da densidade em \eqref{eq:densidade-transformada} é feita
a menos de conjuntos de medida nula. De fato, duas funções não negativas
que coincidem em quase todo ponto produzem as mesmas integrais sobre
qualquer boreliano e, portanto, as mesmas probabilidades. Assim, a
igualdade $\mathbb{P}((U, V) \in B) = \iint_B h$ não distingue duas
densidades que difiram apenas em um conjunto de medida nula. A escolha em
\eqref{eq:densidade-transformada}, que atribui o valor $0$ fora de $S$, é
a representação canônica dessa classe.

\begin{example}[Transformação com Jacobiano e uma Uniforme]\label{ex:transformacao-jacobiano}
Suponha que $X$ e $Y$ tenham função densidade conjunta
\begin{equation}\label{eq:densidade-conjunta-jacobiano}
  f_{X,Y}(x,y) =
  \begin{cases}
    e^{-x-y}, & \text{se } x, y > 0, \\[0.2cm]
    0, & \text{caso contrário},
  \end{cases}
\end{equation}
e defina
\[
  U = X + Y, \qquad V = \frac{X}{X+Y}.
\]
Vamos determinar a densidade conjunta de $(U,V)$ e a densidade marginal de
$V$.

\begin{solution}
  A transformação deste exemplo é $T(x,y) = (u(x,y)\, ,\, v(x,y))$, dada por
  \begin{equation}\label{eq:transformacao-jacobiano-exemplo}
    u(x,y) = x + y, \qquad v(x,y) = \frac{x}{x+y}.
  \end{equation}
  Para usar a fórmula do jacobiano, \Cref{thm:formula-jacobiano}, precisamos
  primeiro verificar que $T$ é uma bijeção de
  $D = \{(x,y) : x, y > 0\}$ em $S = \{(u,v) : u > 0,\ 0 < v < 1\}$.
\begin{figure}[H]
\centering
\begin{tikzpicture}[scale=1.05, >=Stealth]
  % Janela do primeiro quadrante: D se estende para cima e para a direita.
  \fill[teal!15] (0,0) rectangle (3.0,2.8);
  \draw[->, thick] (-0.42,0) -- (3.35,0) node[below] {$x$};
  \draw[->, thick] (0,-0.42) -- (0,3.15) node[left] {$y$};
  % As fronteiras tracejadas não pertencem às regiões.
  \draw[teal, dashed, line width=0.9pt] (0,2.8) -- (0,0) -- (3.0,0);
  \node[below left=2pt] at (0,0) {$0$};
  \node[teal!80!black] at (1.5,1.4) {$D$};

  % Janela da faixa S: a altura 2.1 representa v=1.
  % Aumente ou diminua este valor para ajustar a separação dos gráficos.
  % Com scale=1.05, cada 1 cm aqui corresponde a 1.05 cm no PDF.
  \def\distanciaGraficos{7.2cm}
  \begin{scope}[xshift=\distanciaGraficos]
    \coordinate (setaTchegada) at (0.6,2.9);
    \coordinate (setaTinversaPartida) at (0.6,-0.65);
    \fill[teal!15] (0,0) rectangle (3.0,2.1);
    \draw[->, thick] (-0.42,0) -- (3.35,0) node[below] {$u$};
    \draw[->, thick] (0,-0.42) -- (0,3.15) node[left] {$v$};
    \draw[teal, dashed, line width=0.9pt]
      (3.0,0) -- (0,0) -- (0,2.1) -- (3.0,2.1);
    \node[below left=2pt] at (0,0) {$0$};
    \node[left=3pt] at (0,2.1) {$1$};
    \fill[muted] (0,2.1) circle (1.5pt);
    \node[teal!80!black] at (1.5,1.05) {$S$};
  \end{scope}

  \draw[->, teal!80!black, line width=0.9pt]
    (3.1,1.3) to[out=20, in=160]
    node[midway, above=4pt] {$T$} (6.8,1.3);
  \draw[->, teal!80!black, line width=0.9pt]
    (setaTinversaPartida) to[out=-145, in=-35]
    node[midway, below=4pt] {$T^{-1}$} (2.4,-0.65);
\end{tikzpicture}
\caption{A transformação $T(x,y)=\left( x+y\, , \, \displaystyle\frac{x}{x+y} \right)$.}
\label{fig:dominio-faixa-transformacao}
\end{figure}
  \smallskip

Vamos começar mostrando que a aplicação $T$ leva o conjunto $D$ no conjunto $S$, isto é,
$T(D)\subseteq S$. De fato, para quaisquer $x, y > 0$, temos
\begin{align*}
0< x+y = u
\qquad \text{e}\qquad 
0< v = \frac{x}{x+y}<1,
\end{align*}
Logo $T(D) \subseteq S$.

\smallskip
Nosso próximo passo será mostrar que $T(D)=S$ e que $T$ é bijetiva. 
De fato, dado $(u,v) \in S$, o sistema formado pelas duas equações em
\eqref{eq:transformacao-jacobiano-exemplo} tem solução única
\begin{equation}\label{eq:inversa-jacobiano-exemplo}
x = uv, \qquad y = u(1-v).
\end{equation}
Para verificar que esta afirmação é verdadeira, primeiro observamos que 
\begin{align*}
u=x+y \quad \text{e} \quad v = \frac{x}{x+y} \quad \Longrightarrow \quad x = v(x+y) = vu 
\end{align*}
Substituindo em $u=x+y$, o resultado da última igualdade acima concluímos que
\begin{align*}
y = u - x = u - uv = u(1-v). 
\end{align*}
Como $u > 0$ e $0 < v < 1$, segue que
$x = uv > 0$ e $y = u(1-v) > 0$, ou seja, $(x,y) \in D$. 
Isto mostra que a aplicação 
\begin{align*}
T^{-1}(u,v) = (uv\, ,\, u(1-v))
\end{align*}
envia o conjunto $S$ em $D$. Além do mais, segue das igualdades acima que
$T\circ T^{-1}(u,v)=(u,v)$ e $T^{-1}\circ T(x,y)=(x,y)$ e portanto $T^{-1}$
é de fato a inversa de $T$.  

\medskip
O determinante jacobiano de $T^{-1}$ é o determinante da matriz das
derivadas parciais de $x\equiv x(u,v)$ e $y\equiv y(u,v)$, em relação a $u$ e $v$:
\[
J(u,v)
=
\left|
\begin{array}{>{\displaystyle}c@{\hspace{0.8cm}}>{\displaystyle}c}
\frac{\partial x}{\partial u} & \frac{\partial x}{\partial v} 
\\[0.6cm]
\frac{\partial y}{\partial u} & \frac{\partial y}{\partial v}
\end{array}
\right|
=
\begin{vmatrix}
v\ & u 
\\[0.3cm]
1-v\ & -u
\end{vmatrix}
= -uv - u(1-v) = -u.
\]
Portanto $|J(u,v)| = u$ em todo ponto de $S$, já que $u > 0$.

\medskip
Podemos agora aplicar a fórmula da densidade transformada,
\eqref{eq:densidade-transformada}, segundo a qual
\[
f_{U,V}(u,v)
= f_{X,Y}\bigl(x(u,v), y(u,v)\bigr)\,\bigl|J(u,v)\bigr|.
\]
Para $(u,v) \in S$, as relações \eqref{eq:inversa-jacobiano-exemplo}
dão $x = uv > 0$ e $y = u(1-v) > 0$, de modo que o fator
$e^{-x-y}$ em \eqref{eq:densidade-conjunta-jacobiano} vale
\[
e^{-x-y} = e^{-uv}\,e^{-u(1-v)} = e^{-u}.
\]
Multiplicando pelo jacobiano, obtemos
\begin{equation}\label{eq:densidade-conjunta-uv}
f_{U,V}(u,v) =
\begin{cases}
u\,e^{-u}, & \text{se } u > 0 \text{ e } 0 < v < 1, \\[0.2cm]
0, & \text{caso contrário}.
\end{cases}
\end{equation}

\medskip
Vamos calcular a densidade marginal de $V$. Pela definição de densidade
marginal, integramos \eqref{eq:densidade-conjunta-uv} em relação a $u$.
O integrando é nulo fora de $u > 0$, o que justifica reduzir a integral
a esse intervalo:
\begin{align*}
f_V(v)
&= \int_{-\infty}^{\infty} f_{U,V}(u,v)\,du
\\[0.3cm]
&= \begin{cases}
\displaystyle\int_0^{\infty} u\,e^{-u}\,du, & \text{se } 0 < v < 1, \\[0.4cm]
0, & \text{caso contrário}.
\end{cases}
\end{align*}
A integral que aparece é $\Gamma(2) = 1! = 1$, conforme a definição da
função gama nas notas sobre funções de distribuição. Logo
\begin{equation}\label{eq:marginal-v-uniforme}
f_V(v) =
\begin{cases}
1, & \text{se } 0 < v < 1, \\[0.2cm]
0, & \text{caso contrário},
\end{cases}
\end{equation}
o que mostra que $V$ tem distribuição uniforme no intervalo $(0,1)$.

\medskip
De modo análogo, a marginal de $U$ vem de integrar
\eqref{eq:densidade-conjunta-uv} em relação a $v$, com $v$ percorrendo
$(0,1)$:
\[
f_U(u) = \int_0^1 u\,e^{-u}\,dv = u\,e^{-u}, \qquad u > 0.
\]
Comparando com a densidade gama daquelas notas, com parâmetros de forma
$2$ e taxa $1$, cuja expressão é
$\frac{1}{\Gamma(2)}\,u^{2-1}e^{-u} = u\,e^{-u}$ para $u > 0$,
concluímos que $U$ tem distribuição gama com parâmetros $2$ e $1$.
Além disso, como \eqref{eq:densidade-conjunta-uv} se fatora no produto
da marginal \eqref{eq:marginal-v-uniforme} com a marginal de $U$, isto é,
\[
f_{U,V}(u,v) = f_U(u)\,f_V(v) \qquad \text{para todo } (u,v) \in \mathbb{R}^2,
\]
segue que $U$ e $V$ são independentes. Em particular, reobtemos que
$X+Y$ e $X/(X+Y)$ são independentes, com $X+Y$ gama de parâmetros $2$ e
$1$ e $X/(X+Y)$ uniforme em $(0,1)$.
\end{solution}
\end{example}

\begin{exercise}\label{exercise:diferenca-soma-normais}
  Sejam $X$ e $Y$ variáveis aleatórias independentes, cada uma com
  distribuição normal de média $\mu$ e variância $\sigma^2$. Determine a
  função densidade conjunta do par $(U, V)$, em que
  \[
    U = X - Y
    \qquad \text{e} \qquad
    V = X + Y.
  \]
  As variáveis aleatórias $U$ e $V$ são independentes?
\end{exercise}

\begin{exercise}\label{exercise:distribuicao-dupla-exponencial}
  Sejam $X$ e $Y$ variáveis aleatórias com função densidade conjunta
  \[
    f(x, y) =
    \begin{cases}
      \frac{1}{4}e^{-\frac{1}{2}(x+y)}, & \text{se } x > 0 \text{ e } y > 0, \\[0.3cm]
      0, & \text{caso contrário}.
    \end{cases}
  \]
  Mostre que a função densidade conjunta do par $(U, V)$, em que
  \[
    U = \frac{1}{2}(X - Y)
    \qquad \text{e} \qquad
    V = Y,
  \]
  é dada por
  \[
    f_{U,V}(u, v) =
    \begin{cases}
      \frac{1}{2}e^{-u-v}, & \text{se } (u, v) \in A, \\[0.3cm]
      0, & \text{caso contrário},
    \end{cases}
  \]
  em que $A$ é uma região do plano $(u, v)$ que deve ser determinada.
  Deduza que $U$ tem distribuição exponencial bilateral, com função densidade
  \[
    f_U(u) = \frac{1}{2}e^{-|u|}
    \qquad \text{para todo } u \in \mathbb{R}.
  \]
\end{exercise}

\begin{exercise}\label{exercise:probabilidade-raizes-reais}
  Sejam $B$ e $C$ variáveis aleatórias independentes, cada uma com
  distribuição uniforme no intervalo $[-1,1]$. Calcule a probabilidade da
  equação quadrática aleatória
  \[
    z^2 + Bz + C = 0
  \]
  possuir duas raízes reais, admitindo a possibilidade de uma raiz dupla.
  Use o discriminante $\Delta = B^2 - 4C$ e siga as etapas:
  \begin{enumerate}
    \item Considere a transformação $T(x,y)=(u(x,y)\, ,\, v(x,y))$ induzida por 
    $U=B$ e $V=B^2-4C$. Determine sua
      inversa e a imagem $S$ do quadrado $(-1,1)\times(-1,1)$ no plano $(u,v)$,
      esboçando essa região.
    \item Calcule o jacobiano da transformação inversa. Use que a densidade
      conjunta do vetor aleatório $(B,C)$ é $1/4$ no quadrado para obter a densidade conjunta
      de $(U,V)$ em $S$.
    \item Expresse o evento de a equação ter raízes reais em termos de
      $(U,V)$ e calcule sua probabilidade integrando a densidade transformada
      sobre a parte correspondente de $S$. (Sugestão: note que $|u| < 1$
      implica $u^2-4 < 0$, o que determina o limite inferior de integração
      em $v$.)
    \item Como conferência, calcule diretamente, no plano $(b,c)$, a área da
      região $\{(b,c) \in (-1,1)^2 : c \le b^2/4\}$ e conclua que a
      probabilidade, dada pela razão entre essa área e a área $4$ do
      quadrado, é $13/24$.
  \end{enumerate}
\end{exercise}


\bibliographystyle{plain}
\bibliography{../arquivos-aux/referencias}
\end{document}
